\documentclass[11pt,a4paper]{article}
\usepackage[T1]{fontenc}
\usepackage[utf8]{inputenc}
\usepackage{newtxtext,newtxmath}
\usepackage[a4paper,margin=33mm]{geometry}
\usepackage{microtype}
\usepackage{graphicx}
\usepackage{xcolor}
\usepackage{mdframed}
\let\openbox\relax
\usepackage{amsmath,amsthm,amssymb,mathtools}
\usepackage{needspace}
\usepackage[explicit]{titlesec}
\usepackage{titling}
\usepackage[
  pdfusetitle,
  hidelinks,
  bookmarksopen=true,
  bookmarksnumbered=true
]{hyperref}

\linespread{0.95}
\setlength{\parindent}{0pt}
\setlength{\parskip}{5.5pt}
\allowdisplaybreaks
\numberwithin{equation}{section}

\DeclareMathOperator{\Tr}{Tr}
\DeclareMathOperator{\rank}{rank}
\DeclareMathOperator{\diag}{diag}
\DeclareMathOperator{\spanop}{span}
\newcommand{\R}{\mathbf R}

\titleformat{\section}
{\normalfont\normalsize\bfseries\raggedright}
{\thesection}{0.75em}{\begin{minipage}[t]{0.9\textwidth}#1\end{minipage}}
[\vspace{0.12em}\titlerule]

\newtheoremstyle{thmcompact}
{0.5ex}{0.5ex}
{\itshape}
{0pt}
{\bfseries}
{.}
{0.6em}
{\thmname{#1}\ \thmnumber{#2}\ \thmnote{(\textit{#3})}}

\theoremstyle{thmcompact}
\newtheorem{theorem}{Theorem}[section]
\newtheorem{proposition}[theorem]{Proposition}
\newtheorem{corollary}[theorem]{Corollary}
\newtheorem{definition}[theorem]{Definition}
\newtheorem{lemma}[theorem]{Lemma}

\theoremstyle{remark}
\newtheorem*{remarkplain}{Remark}

\mdfdefinestyle{thmframe}{%
leftline=true, rightline=false, topline=false, bottomline=false,
linewidth=0.8pt, linecolor=black,
innerleftmargin=16pt, innerrightmargin=8pt,
innertopmargin=6pt, innerbottommargin=6pt,
skipabove=6pt, skipbelow=6pt,
nobreak=true
}
\surroundwithmdframed[style=thmframe]{theorem}
\surroundwithmdframed[style=thmframe]{proposition}
\surroundwithmdframed[style=thmframe]{corollary}
\surroundwithmdframed[style=thmframe]{definition}
\surroundwithmdframed[style=thmframe]{lemma}

\newmdenv[
linewidth=0.6pt,
topline=false,
bottomline=false,
rightline=false,
skipabove=\baselineskip,
skipbelow=\baselineskip,
backgroundcolor=gray!3,
leftmargin=0pt,
innerleftmargin=8pt,
innerrightmargin=6pt,
frametitleaboveskip=0.4em,
frametitlebelowskip=0.2em,
frametitlefont=\bfseries\small,
]{remarkbox}
\newenvironment{remark}[1][]%
{\begin{remarkbox}\begin{remarkplain}[#1]\normalfont}
{\end{remarkplain}\end{remarkbox}}

\renewcommand{\qedsymbol}{}
\newcommand{\thmneed}{\Needspace{6\baselineskip}}
\AtBeginEnvironment{theorem}{\thmneed}
\AtBeginEnvironment{proposition}{\thmneed}
\AtBeginEnvironment{corollary}{\thmneed}
\AtBeginEnvironment{definition}{\thmneed}
\AtBeginEnvironment{lemma}{\thmneed}

\title{\textbf{Claude Notes Audit Addendum}\\
\large{Differential-correlation calculus, noncommuting closure, and local-geometry ensembles}}
\author{J.\,R.~Dunkley}
\date{11 April 2026}

\hypersetup{
  pdftitle={Claude Notes Audit Addendum},
  pdfauthor={J. R. Dunkley},
  pdfsubject={Audit of external Claude research notes for observer geometry},
  pdfkeywords={observation geometry, Gaussian covariance, branch selection, closure, Fisher geometry}
}

\begin{document}

\maketitle

\begin{abstract}
The notes in \texttt{papers/claude\_notes} are useful, but only after triage.
Their strongest exact content is covariance-level observer geometry already
native to the project: visible precision is marginal precision, the determinant
clock is a log-det covariance correlation, tower laws compose observers, and
closure leakage is exact for supplied quadratic families. The new value is more
specific. A rank-one differential-correlation perturbation gives an exact
one-channel hidden-gap update only in a white or aligned background; in a
coloured background the update is generally two-rank. Generic noncommuting
perturbation families need not admit exact intermediate-rank closure observers.
Branch and observer selection should therefore be stated with residual margins,
not promoted to generic non-Gaussian law selection. The broadest opportunity is
a formal ensemble layer over local quadratic geometries, but that is a higher
order/statistical extension, not a replacement for full-law cumulants.
\end{abstract}

\section{Triage judgement}

The Claude notes should be treated as a hypothesis source. They should not be
merged into the core papers as application claims. Their black-hole, neural, and
turbulence notes mostly police their own scope: they operate at Gaussian
covariance level or at exact local quadratic/Fisher level. That part is
consistent with the current project doctrine.

What survives as project-native mathematics is narrower and cleaner:
\[
  \Phi_C(H)=(CH^{-1}C^{\mathsf T})^{-1},
  \qquad
  \tau=\log\det T-\log\det\Phi,
  \qquad
  \mathcal L(\Pi)=\frac12\sum_a\|[D_a,\Pi]\|_F^2.
\]
These are exact quadratic objects. They become full-law statements in Gaussian
or other explicitly solved sectors, such as the fixed-precision affine-hidden
sector. The notes do not justify a generic non-Gaussian branch selector.

\section{Differential-correlation sector}

The neural note contains a useful but overbroad phrase: a differential
correlation creates exactly one leakage channel aligned with the signal
direction. The statement is exact in the tested model, but its theorem-level
scope is smaller.

\begin{proposition}[White rank-one covariance sector]
Let
\[
  \Sigma_\varepsilon=I+\varepsilon ff^{\mathsf T},
  \qquad
  f=(a,b)\in\R^m\oplus\R^{n-m},
  \qquad
  \varepsilon>0.
\]
Let \(H_\varepsilon=\Sigma_\varepsilon^{-1}\), \(T=(H_\varepsilon)_{VV}\),
and \(\Phi=(\Sigma_{\varepsilon,VV})^{-1}\). Then
\[
  T-\Phi
  =
  \frac{\varepsilon^2\|b\|^2}
       {(1+\varepsilon\|a\|^2)(1+\varepsilon\|f\|^2)}
       aa^{\mathsf T}.
\]
Thus the hidden gap is rank one if \(a\ne0\) and \(b\ne0\), and its visible
channel is the visible signal direction \(a\).
\end{proposition}

\begin{proof}
Sherman--Morrison gives
\[
  H_\varepsilon
  =
  I-\frac{\varepsilon}{1+\varepsilon\|f\|^2}ff^{\mathsf T},
\]
so
\[
  T
  =
  I_m-\frac{\varepsilon}{1+\varepsilon(\|a\|^2+\|b\|^2)}aa^{\mathsf T}.
\]
Also
\[
  \Phi
  =
  (I_m+\varepsilon aa^{\mathsf T})^{-1}
  =
  I_m-\frac{\varepsilon}{1+\varepsilon\|a\|^2}aa^{\mathsf T}.
\]
Subtracting yields the displayed expression.
\end{proof}

\begin{proposition}[Coloured background is generally two-rank]
Let
\[
  \Sigma_\varepsilon=\Sigma_0+\varepsilon ff^{\mathsf T},
  \qquad H_0=\Sigma_0^{-1},
  \qquad \Phi_0=(\Sigma_{0,VV})^{-1}.
\]
Write \(u_V=(H_0f)_V\) and \(w_V=\Phi_0 f_V\). Then the hidden-gap increment
\[
  \Delta R:=(T_\varepsilon-\Phi_\varepsilon)-(T_0-\Phi_0)
\]
equals
\[
  \Delta R
  =
  \frac{\varepsilon}{1+\varepsilon f_V^{\mathsf T}\Phi_0f_V}w_Vw_V^{\mathsf T}
  -
  \frac{\varepsilon}{1+\varepsilon f^{\mathsf T}H_0f}u_Vu_V^{\mathsf T}.
\]
Consequently \(\rank(\Delta R)\le2\), and the one-channel conclusion requires
the two visible directions \(u_V\) and \(w_V\) to be collinear.
\end{proposition}

\begin{proof}
Apply Sherman--Morrison once to the full covariance, which gives the update of
\(T_\varepsilon=(H_\varepsilon)_{VV}\), and once to the visible covariance
\((\Sigma_\varepsilon)_{VV}=(\Sigma_0)_{VV}+\varepsilon f_Vf_V^{\mathsf T}\),
which gives the update of \(\Phi_\varepsilon\). Subtract.
\end{proof}

The audit numerics found a coloured-background example with
\[
  \rank(\Delta R)=2,\qquad
  \|\Delta R-\Delta R_{\rm formula}\|_{\max}=1.18\cdot10^{-16},
  \qquad
  |\langle u_V,w_V\rangle|/(\|u_V\|\|w_V\|)=0.809.
\]
Thus the Claude claim is not wrong in its toy sector, but it should be promoted
as a rank-one solvable subcase, not as a generic law of differential
correlations.

\section{Noncommuting closure}

The neural note also reports exact \(\eta=0\) closure-adapted observers at
several ranks. The verifier confirms this, but the construction uses a
commuting family built in an eigenbasis. That is a legitimate exact mode, not a
generic noncommuting result.

\begin{proposition}[Exact closure can fail generically]
Let \(D_1=\diag(1,\ldots,n)\) and let \(D_2\in\mathrm{Sym}(n)\). For a fixed
rank \(m\), an exact closure subspace for both \(D_1\) and \(D_2\) exists only
if some coordinate \(m\)-plane has zero \(D_2\) cross-block. For a generic
\(D_2\), this condition fails for every intermediate rank.
\end{proposition}

\begin{proof}
Since \(D_1\) has simple spectrum, every \(D_1\)-invariant subspace is spanned
by a subset of the coordinate eigenvectors. Exact closure for \(D_2\) then
requires \(D_2(U)\subseteq U\), equivalently the \(U^\perp\times U\) block of
\(D_2\) is zero. For each fixed coordinate subset this imposes polynomial
equalities on the entries of \(D_2\). A finite union of such proper algebraic
sets has empty interior and measure zero. Hence generic \(D_2\) gives no exact
common intermediate-rank closure subspace.
\end{proof}

In the audit script with \(n=6\), \(m=2\), the commutator norm was \(6.80\) and
the smallest coordinate cross-block certificate was \(1.113\). A random
Grassmann search found best leakage \(1.376\), consistent with the compactness
consequence that the infimum is positive once exact common invariance is
excluded.

\section{Robust branch and observer claims}

The durable upgrade from the Claude notes is not ``the best observer is exact
closure''. The durable upgrade is the margin doctrine. Suppose two quadratic
branch scores satisfy \(q_2-q_1=\gamma>0\), and suppose unknown law-level or
model-level residuals obey \(|r_i|\le R\). Then the ordering is forced only when
\(\gamma>2R\). The same logic applies to selector Hessians: if the quadratic
Hessian gap is \(\gamma\) and the residual Hessian operator norm is less than
\(\gamma\), strict local stability persists.

The audit script used \(\gamma=0.08\). With \(R=0.03\), no reversal occurred in
50,000 residual draws and none is possible by the theorem. With \(R=0.05\), the
condition fails; a constructive reversal has residuals \(r_1=0.05\),
\(r_2=-0.05\), giving full gap \(-0.02\), and random draws produced 953
reversals out of 50,000.

\section{Local-geometry ensembles}

The turbulence note's useful conceptual move is not that intermittency has been
derived. It has not. The useful move is to name a possible higher layer:
instead of one Hessian \(H\), use a distribution \(\mu(dH)\) of local Hessians
or Fisher objects. Each fibre \(H\) has exact quadratic observer geometry, while
the distribution over \(\Phi_C(H)\), \(\Lambda(H)\), and \(\tau(H)\) carries
heterogeneity and tail information that a single averaged Hessian can lose.

The audit numerics constructed two local-coupling ensembles with matched mean
coupling \(0.13315\). The clock at the mean coupling was \(0.02363\), while the
mean clocks were \(0.03734\) and \(0.03455\). The 99th percentile clocks were
\(0.38274\) and \(0.48870\). Thus
\[
  \tau(\mathbf E H_{\rm local}) \ne \mathbf E\,\tau(H_{\rm local}),
\]
and tail-sensitive observational claims cannot be reduced to one representative
quadratic object. This supports an ensemble-of-local-geometries extension, not a
generic full non-Gaussian law theorem.

\section{Application-note triage}

\textbf{Black-hole note.} Keep only the covariance-level dictionary and the
branch-minimised observer-prescription analogy. Do not identify the determinant
clock with von Neumann entropy except under separately stated Gaussian/symplectic
conditions. The area term, replica topology, and matrix-element encoding remain
outside the project.

\textbf{Neural note.} Keep the Fisher quadratic bridge, Sherman--Morrison
precision perturbation, exact FI ceiling in the linear Gaussian model, and the
rank-one white differential-correlation sector. Amend the general leakage-channel
language to the rank-one/two-rank calculus above. Keep closure-adapted
\(\eta=0\) only in commuting or explicitly invariant sectors.

\textbf{Turbulence note.} Keep the statement that filters are observation maps
and that cross-scale covariance coupling has exact hidden-load geometry. Treat
triadic cascade exponents, spectral concentration, anomalous dissipation, and
intermittency claims as model observations or structural resemblances unless a
separate dynamical proof is supplied. The strong next theory idea is the
local-geometry ensemble, not a claimed derivation of anomalous exponents.

\section{Module and theory implications}

No module source change is licensed by these notes. The module already supports
the exact quadratic primitives. Sensible future additions would be theorem-local:

\begin{itemize}
\item a rank-one covariance perturbation helper that reports the white/aligned
one-channel case and the general two-rank coloured update;
\item a closure irreducibility diagnostic for simple-spectrum test families,
returning certificates that exact intermediate-rank closure is impossible;
\item a residual-margin wrapper for branch and observer selection;
\item an ensemble evaluator that maps samples of local Hessians to distributions
of \(\Phi_C\), \(\Lambda\), and \(\tau\), explicitly labelled as an ensemble
quadratic layer.
\end{itemize}

The broader judgement is unchanged. The project is past Gaussian toy territory
because its exact local quadratic geometry, fixed-precision affine-hidden sector,
and rank-one perturbation sectors are genuinely structural. It is not past the
need for law-level hypotheses. Higher cumulants, fibre-volume terms, remote
mixture mass, and local-geometry distributions are real next frontiers, not
cosmetic relabellings.

\end{document}
